%==============================================================================
\section{madGLP Communication Model}
\label{sec:madglp-comm}
%==============================================================================
This section summarizes the aspects of madGLP that are relevant to the networking layer. The full specification is in~\cite{shapiro2026implementing}.
%------------------------------------------------------------------------------
\subsection{Global Links and Messages}
\label{sec:global-links}
%------------------------------------------------------------------------------
In madGLP, inter-agent communication is realized through \emph{global links}. A global link connects a local variable pair $(X_p, X_p?)$ at agent $p$ to a local variable pair $(X_q, X_q?)$ at agent $q$. When the writer at one end is assigned, the value is sent to the other end via a message.
All inter-agent messages have the form:
$$G := T^\uparrow$$
where $G$ is a \emph{global variable name} identifying the target endpoint, and $T^\uparrow$ is a \emph{globalized term}---a term in which local variables have been replaced by global names for transmission.
A global variable name has the form $\_w(p, i)$ or $\_r(p, i)$, where $p$ is an agent identifier and $i \in \mathbb{N}$ is an index. The name encodes both the originating agent and the role (writer or reader) of the variable that was globalized.
%------------------------------------------------------------------------------
\subsection{Message Types}
\label{sec:msg-types}
%------------------------------------------------------------------------------
There are two categories of messages, distinguished by the index in the global name:
\textbf{Regular link messages} (index $i > 0$): A point-to-point assignment sent from the agent where a \code{global\_send} goal fires to the agent that holds the corresponding entry in its global writers table. The message is $G := T^\uparrow$ where $G$ is either $\_w(p,i)$ or $\_r(p,i)$. Upon receipt, the target agent looks up the entry, assigns the local writer, and removes the entry.
\textbf{Cold-call messages} (index $i = 0$): A message to an agent's well-known index-0 serializer, used to initiate communication between agents that share no prior variables. The message is $\_w(q,0) := [T^\uparrow \mid \_w(q,0)]$, wrapping the term in a list cell. Upon receipt, the target agent appends the term to its network input stream and updates (but does not remove) the serializer entry. Multiple agents may cold-call the same target concurrently.
%------------------------------------------------------------------------------
\subsection{Message Flow}
\label{sec:msg-flow}
%------------------------------------------------------------------------------
The lifecycle of a madGLP message involves three steps, corresponding to three madGLP transactions:
\textbf{Step 1 --- Reduce:} A writer $X$ is assigned at agent $p$. The paired reader $X?$ becomes known, triggering a \code{global\_send(X?, G, Q)} goal. The goal's reduction globalizes $X?$'s value and adds message $(G := T^\uparrow, Q)$ to $p$'s outgoing message set $M_p$.
\textbf{Step 2 --- Send:} The message $(G := T^\uparrow, Q)$ is removed from $M_p$ and placed in the communication channel to $Q$. \emph{This is the step that the networking layer implements.}
\textbf{Step 3 --- Receive:} The message arrives at agent $Q$. The runtime looks up $G$ in $Q$'s global writers table, localizes $T^\uparrow$ (creating fresh local variable pairs for any global names it contains), assigns the local writer, and reactivates any suspended goals.
The networking layer is responsible only for Step~2: delivering a serialized message from agent $p$ to agent $Q$. Steps~1 and~3 are handled entirely within the GLP runtime on each device.
